\section{Proofs for call-by-name PCF adequacy}
\label{annex:p10-operational-adequacy-domain-models}

\begin{annexitem}{annex:p10.oa.operational}{thm:p10-oa-operational}{Typed substitution and complete evaluation correspondence}
\begin{proof}
For typed substitution, induct on the typing derivation. A variable is
either replaced by its same-typed supplied term or remains a declaration
in the residual context. Constants are unaffected; each unary constructor
uses the induction hypothesis on its argument; conditional uses it on
guard and both branches; application uses it on both typed subterms.
For abstraction, first choose the bound name fresh for the substitution's
range; extend the residual context by its original type and apply the
induction hypothesis to the body. The abstraction rule then restores
the claimed arrow type. For \(\mu x:t.M\), choose \(x\) fresh in
the same way; the body has type \(t\) under its \(x:t\) declaration,
so substitution and the recursion rule give type \(t\) again.
This also proves simultaneous substitution by induction, retaining each
binder's annotation and not substituting into newly inserted terms.

Each primitive contraction preserves typing: the beta and recursive
cases use the just-proved substitution theorem; predecessor/zero-test
use the numeral types; conditionals retain the branch type. Induction
on the evaluation-context grammar lifts that property to every step.
Canonical forms are immediate from the typed value grammar: a numerical
value is a numeral, a boolean value is true or false, and an arrow value
is an abstraction. For progress, induct on the syntax of a closed typed
term. Variables are impossible. Constants and abstractions are values;
successor steps its argument unless that argument is a numeral, when it
is a value. Predecessor and zero-test step their argument or use the
appropriate numeral contraction. Application steps only its operator
until it is an abstraction, when beta applies. Conditional steps only
its guard until it is a boolean, then selects exactly one branch.
Recursion always unfolds. These alternatives are exclusive; the induction
hypothesis on the unique active subterm proves determinism. Values have
no primitive contraction or active substep.

We next prove a one-step evaluation equivalence, before using the
small/big-step correspondence. At a beta redex, inversion of NameApp
and the abstraction axiom says precisely that its substituted body
evaluates to the requested value. At a recursion redex, RecEval is
exactly the equivalence with its unfolded body. At the two conditional
redexes, the guard axiom and branch rule give the selected branch.
At predecessor and zero-test redexes, the numeral evaluates to itself,
by induction on the numeral, and inversion of the two arithmetic rules
gives exactly the displayed contractum. Thus primitive \(M\to N\)
preserves and reflects \(\Downarrow V\), for every value \(V\).
To lift through \(E\), induct on its grammar. Succ, Pred and Test
replace exactly the argument-evaluation premise. An operator context
\(E[-]L\) replaces the premise evaluating the operator to the same
abstraction, leaving the substituted-body premise identical. A guard
context replaces its boolean-evaluation premise and leaves the selected
branch identical. This proves the equivalence for every small step.

If \(M\Downarrow V\), induction on its finite derivation constructs
\(M\to^*V\). A successor evaluates its argument under Succ contexts;
predecessor/test then perform their numeral contraction. Application
first evaluates its operator under the operator context, contracts beta,
and follows the body derivation. Conditional first evaluates its guard,
contracts the chosen branch and follows that derivation. Recursion
unfolds once and follows its premise. Constant and abstraction cases
use zero steps. Conversely every value evaluates to itself, by numeral
induction and the base/abstraction axioms. Along a finite reduction
sequence ending in \(V\), repeatedly apply the one-step equivalence
backwards to obtain \(M\Downarrow V\).
Determinism and irreducibility of values imply uniqueness of that result.
Finally, preservation and progress allow a next step at every nonvalue
stage of a closed typed computation. If there is no finite result,
iteration of this unique next step gives an infinite sequence. If there
is a finite result, its deterministic sequence reaches an irreducible
value and cannot also continue infinitely. This proves the stated
divergence characterization, not an algorithm deciding divergence.
\end{proof}
\end{annexitem}

\begin{annexitem}{annex:p10.oa.semantic-soundness}{thm:p10-oa-semantic-soundness}{Continuity, substitution and the sound direction}
\begin{proof}
The flat ground carriers are pointed bounded-complete algebraic domains.
Theorem~\ref{thm:p10-domain-algebraic-functions}, iterated on the type
grammar, supplies the same properties for every full continuous function
space. Finite products form environment dcpo's.
Interpret terms by induction. Projections/constants are continuous.
The ground operations are continuous: a directed set in a flat carrier
either consists only of bottom or contains a unique nonbottom value,
which is its supremum. Their strict definitions preserve those suprema.
For a directed family of guard/branch triples, a nonbottom guard in the
supremum appears in some triple and then is constant on a cofinal upper
tail. On that tail the selected branch's directed supremum is exactly
the conditional result. If the supremum guard is bottom, all guard
values and all results are bottom. This proves joint continuity of the
conditional for every result type, including arrows.
Application and abstraction use continuous evaluation and currying from
Theorem~\ref{thm:p10-domain-function-space}. For recursion the body
induces a continuous map from environments into \([D_t\to D_t]\);
the continuous least-fixed-point operator of
Theorem~\ref{thm:p10-domain-least-fixpoint} composes with it. Thus all
terms have jointly continuous meanings and the recursion clause exists.

For semantic substitution, induct on term construction. Variables and
constants give the asserted equality immediately, and operations and
application combine the equalities for their arguments. In each binder,
choose the bound name fresh for all substituends. For abstraction, evaluate
the two functions at arbitrary \(d\); the body induction hypothesis
under the updated environment proves equal values. For recursion,
the body induction hypothesis under \(x\mapsto d\), for every \(d\),
makes the two defining continuous endomaps equal. Their least fixed
points are therefore equal. The same fresh-name argument proves semantic
bound-renaming invariance. This proves the simultaneous substitution
formula, including recursive subterms, rather than importing an STLC-only
substitution assertion for the enlarged syntax.

Beta soundness is evaluation/currying plus substitution. The recursive
contraction is sound because its least fixed point \(a\) satisfies
\(a=\llbracket M\rrbracket_{\rho[x\mapsto a]}\); substitution identifies
that expression with the unfolded term's denotation. Numeral and
conditional contractions have exactly the stated semantic values,
including predecessor at zero. Compositionality lifts equality through
every evaluation context. Consequently every finite sequence preserves
denotation, and the operational correspondence gives soundness for
\(M\Downarrow V\). No argument here yet reflects nonbottom meaning
back into a finite operational derivation.
\end{proof}
\end{annexitem}

\begin{annexitem}{annex:p10.oa.relation-closure}{thm:p10-oa-relation-closure}{The admissible relation is closed at every type}
\begin{proof}
Induct on type. At ground type bottom is related by definition.
For a directed family of related elements, a nonbottom supremum is a
compact flat value and must occur among them. Its relatedness already
gives the required evaluation. If the supremum is bottom, relatedness
is automatic. At \(s\to t\), the bottom function sends every related
argument to bottom, related at \(t\) by induction. For a directed
family \((f_i)\) of functions related to \(M\), fix any
\(a\mathrel{\mathcal R_s}P\). The directed family \(f_i(a)\) is
related to \(M P\). The induction hypothesis at \(t\) relates its
supremum to \(M P\); pointwise function suprema give the desired
arrow clause. Thus every fixed-term fiber contains bottom and is closed
under all directed suprema, not just the chain used later.

For a one-step conversion at ground type, relatedness is either automatic
bottom or precisely the evaluation equivalence established in
Theorem~\ref{thm:p10-oa-operational}. At an arrow type, for every
related argument \(a,P\), operator-context propagation gives
\(M P\to N P\). The induction hypothesis at the result type equates
their relatedness with \(f(a)\). Quantifying over all such arguments
proves both directions of the arrow equivalence. Iteration gives the
claim for \(\to^*\). This is an induction on types, not an unproved
conversion axiom imposed on the relation.
\end{proof}
\end{annexitem}

\begin{annexitem}{annex:p10.oa.fundamental}{thm:p10-oa-fundamental}{Every syntax clause and the complete recursion induction}
\begin{proof}
Induct on the typing derivation, proving the assertion for all related
environments and closed substitutions. A variable uses the environment
premise; zero and the booleans evaluate to their denoted values. For
application, the two induction hypotheses give a related function and
related argument; the arrow clause gives their application directly.
For \(\lambda x:s.L\), choose \(x\) fresh. Given any closed
\(P:s\) and \(a\mathrel{\mathcal R_s}P\), extend the related
environment/substitution by \(x\mapsto a,P\). The body induction
hypothesis relates \(\llbracket L\rrbracket_{\rho[x\mapsto a]}\)
to \(L[\sigma][P/x]\). Beta conversion and the relation-closure
theorem relate that same element to \((\lambda x:s.L[\sigma])P\).
This is exactly the required arrow clause.

For successor/predecessor/test, if the input meaning is bottom, strictness
makes the output bottom, related to every corresponding term. Otherwise
the ground induction hypothesis supplies its numeral evaluation.
The registered arithmetic rule gives the required result, including
\(\operatorname{pred}(0)\Downarrow0\) and the two zero-test outcomes.
For conditional, a bottom guard gives bottom at its result type. A
nonbottom guard is related to its boolean result, so the operational
correspondence supplies a finite reduction of the guard to that boolean.
Guard-context propagation and the branch contraction reduce the whole
conditional to the selected substituted branch. Its induction hypothesis
and all-type conversion closure establish the desired relatedness.
No premise about the unselected branch's termination is used.

For \(\mu x:t.L\), let \(Q=\mu x:t.L[\sigma]\), with \(x\)
fresh and excluded from \(\sigma\), and define
\[
 F(d)=\llbracket L\rrbracket_{\rho[x\mapsto d]},\qquad
 a_0=\bot_t,\quad a_{n+1}=F(a_n).
\]
We prove \(a_n\mathrel{\mathcal R_t}Q\) by induction on \(n\).
Bottom closure proves the base. For the step, the induction hypothesis
makes \(x\mapsto a_n,Q\), together with the original assignments,
a related environment/substitution for the smaller body derivation.
The typing induction hypothesis for that body therefore gives
\(F(a_n)\mathrel{\mathcal R_t}L[\sigma][Q/x]\).
The recursion rule gives \(Q\to L[\sigma][Q/x]\), so backwards
conversion closure relates \(a_{n+1}\) to \(Q\).
This is not circular: the body theorem is the structural induction
hypothesis, while relatedness of the substitution term \(Q\) to each
finite \(a_n\) is the separate numerical induction just established.
Continuity makes \((a_n)\) increasing, its supremum is the proved
least fixed point, and directed closure of the fixed-term fiber gives
\(\bigsqcup_n a_n\mathrel{\mathcal R_t}Q\).
That is exactly the recursion denotation related to its full recursive
syntax. These cases exhaust the registered typing rules.
\end{proof}
\end{annexitem}

\begin{annexitem}{annex:p10.oa.adequacy}{thm:p10-oa-adequacy}{Neither observation direction is left implicit}
\begin{proof}
For an empty context the fundamental lemma yields
\(\llbracket M\rrbracket\mathrel{\mathcal R_g}M\).
If the denotation equals that of a ground value \(V\), it is nonbottom
and the ground relation clause gives \(M\Downarrow V\). Conversely,
operational evaluation soundness gives that equality whenever
\(M\Downarrow V\). The flat ground interpretation distinguishes
every numeral and the two booleans, so the indicated value is exactly
the observed one, not merely some terminating value.
If a program has bottom denotation, soundness and nonbottom value
interpretations prohibit any evaluation. If it has no evaluation, its
denotation cannot be any nonbottom ground point by the converse direction;
the complete list of flat ground points therefore makes it bottom.
The operational theorem equates this nontermination with the infinite
typed evaluation sequence. No higher-type value test enters the argument.
\end{proof}
\end{annexitem}

\begin{annexitem}{annex:p10.oa.contextual-soundness}{thm:p10-oa-contextual-soundness}{All typed closing contexts, including binding contexts}
\begin{proof}
Suppose the meanings of \(M\) are below those of \(N\) in every
environment of their common context. Induct on a typed one-hole program
context. At the hole this is the assumption. At a ground operation,
application or conditional, monotonicity of the interpreted constructor
preserves the inequality, keeping all other subterms equal. Under
abstraction, the hypothesis holds for every updated environment, so it
gives pointwise order of the two functions. Under a recursion binder,
the two body functions are pointwise ordered; monotonicity of iteration
from bottom and of directed suprema gives order of their least fixed
points. Thus even when context binders intentionally close variables of
the inserted terms, the all-environments assumption applies to the
environments induced by those binders. Capture-avoiding substitution is
not mistakenly used to remove that specified binding.

For a closing ground context obtain
\(\llbracket C[M]\rrbracket\sqsubseteq\llbracket C[N]\rrbracket\).
If \(C[M]\Downarrow V\), adequacy makes the left side the nonbottom
flat value of \(V\). In a flat carrier the only point above that value
is itself, so the right side equals it. The other adequacy direction
gives \(C[N]\Downarrow V\). This proves contextual approximation.
Equality of meanings yields both orders and hence both contextual
approximations. The proof supplies no order or equality in the reverse
direction from quantification over syntactic contexts alone.
\end{proof}
\end{annexitem}

\begin{annexitem}{annex:p10.oa.strategy-boundaries}{thm:p10-oa-strategy-boundaries}{The exact observation and strategy countercases}
\begin{proof}
Unfolding \(\Omega_t\) yields itself. Its unique next-step sequence
never reaches a value, so it has no finite evaluation, and its denotation
is the least fixed point of identity, namely \(\bot_t\).
The call-by-name application contracts directly to zero and evaluates
to zero without inspecting that argument.
For the comparison call-by-value relation, replace NameApp by the rule
requiring both an abstraction result for the operator and a value result
for the argument, then evaluate the value-substituted body; retain the
other natural rules. Any derivation of the example's result would require
a finite derivation of \(\Omega_{\mathsf{num}}\)'s value. No such
derivation exists: its only possible last rule has the identical
judgment as a strictly smaller premise, impossible in a finite tree.
Nevertheless its unchanged call-by-name function interpretation is the
constant-zero function applied to bottom, hence zero. This explicitly
refutes reuse of this interpretation for the modified strategy.

The abstraction \(\lambda x:\mathsf{num}.\Omega_{\mathsf{num}}\)
evaluates to itself and denotes the constant-bottom function. The
least fixed point of identity on the arrow domain is also its bottom
function, which is the denotation of \(\Omega_{\mathsf{num}\to\mathsf{num}}\).
The latter does not evaluate, by the same unfolding argument. Thus the
higher-type denotation alone cannot distinguish termination to a weak
function value from divergence. The adequacy theorem observes only the
ground results reached by program contexts, as its language definition
requires.
\end{proof}
\end{annexitem}
\begin{annexitem}{annex:p10.oa.countdown}{thm:p10-oa-countdown}{A recursive call-by-name derivation and the semantic chain}
\begin{proof}
The recursion typing rule and the conditional/arithmetic rules give
\(K:\mathsf{num}\to\mathsf{num}\). Unfolding once produces the
closed abstraction whose body is the displayed conditional with each
free \(f\) replaced by \(K\); RecEval and AbsValue give that operator
evaluation without evaluating its body.
We prove the assertion for every closed \(P\Downarrow\bar n\) by
induction on \(n\), with this universal quantifier retained in the
induction hypothesis. At zero, ZeroTest evaluates the guard to true and
the selected branch to zero; NameApp therefore gives \(K P\Downarrow0\).
For \(P\Downarrow\operatorname{succ}(\bar n)\), ZeroTest gives false
and PredEval gives \(\operatorname{pred}(P)\Downarrow\bar n\).
The induction hypothesis, applied to that unevaluated expression as a
new argument, gives \(K(\operatorname{pred}(P))\Downarrow\bar n\).
SuccEval and FalseBranch evaluate the body to
\(\operatorname{succ}(\bar n)\), and NameApp gives the asserted
application result. This constructs a finite derivation; it does not
illegally reduce an argument before a call-by-name beta step.

For the semantic chain, let \(F\) be the interpreted displayed body
as an endomap of \([\mathbb N_{\bot}\to\mathbb N_{\bot}]\).
The guard is strict, so \(F(h)(\bot)=\bot\);
\(F(h)(0)=0\), and \(F(h)(n+1)=\operatorname{succ}(h(n))\).
The bottom function has no nonbottom output. If its \(i\)-th iterate
has value \(n\) exactly for \(n<i\), then the next iterate first
gives zero at zero and gives \(n+1\) exactly when \(n<i\). Thus
the \((i+1)\)-st iterate has exactly the claimed support. At each
fixed numeral \(n\), iterate \(n+1\) already gives \(n\), so the
pointwise supremum is identity on all numerals and is bottom at bottom.
The fixed-point theorem identifies this supremum with the denotation
of \(K\). Its application to bottom is bottom; ground adequacy and
the divergence criterion give \(K\Omega_{\mathsf{num}}\Uparrow\).
This explains both the terminating input and the demanded divergent
input in the same specified language/model.
\end{proof}
\end{annexitem}
\begin{annexitem}{annex:p10.oa.sequential-relations}{thm:p10-oa-sequential-relations}{All closure and constructor cases of the semantic relations}
\begin{proof}
Induct on type for the properties of \(Q_t\). At ground type the
first/second-bottom property is in its definition. If a directed supremum
has first and second components nonbottom, flat compactness and
directedness give one triple where both have already appeared. Its
membership forces the third value equal to both; every upper triple
has the same three values. Thus the supremum is related. Otherwise one
of its first two components is bottom and it is related automatically.
At an arrow type, a bottom first or second function gives such a result
on every related argument triple; use induction at the result type.
For directed closure evaluate a directed family of function triples at
a fixed related argument triple, use closure of the result relation,
and use pointwise function suprema. Constantly-bottom tuples are included.
In particular each \(A_t\) contains bottom and is directed-supremum closed.

The ternary semantic fundamental lemma is a typing induction. Variables
use the related environment premise; ground constants give related
diagonals. Strict ground unary operations preserve the relation: if
their first two outputs are nonbottom, both input components are
nonbottom, so all three inputs and outputs are equal. Application uses
the lifted relation and abstraction the body hypothesis on every related
argument triple. For conditional, a bottom first/second guard makes that
output component bottom, invoking the bottom-coordinate result. Otherwise
all guards are the same boolean and select the same related branch.
For recursion, the bottom tuple is related; the body induction hypothesis
preserves related triples of approximants, and directed closure relates
the three least fixed points. These cases exhaust the syntax. Hence a
diagonal environment of members of \(A_t\) gives a diagonal related
result. Closed-term membership in \(A_t\) follows in the empty context.

For \(B_t\), induct on type to prove bottom-pair and directed closure.
Equality at ground type is closed. At arrows the endpoints of a directed
supremum remain in \(A_{s\to t}\); evaluation at every fixed related
argument pair yields a directed \(B_t\)-related family. Closure at
the result type and pointwise suprema prove the arrow clause. Bottom
functions give a bottom output pair on every input. This proves closure
without asserting that \(A_t\) is a full domain or every member definable.

For the binary fundamental lemma, related environment endpoints lie in
the relevant \(A_t\)'s. The ternary lemma separately ensures membership
of every interpreted endpoint in \(A_t\). Variables/constants are
binary-related. Unary ground operations send equal inputs to equal
outputs. Related guards are equal: at bottom both outputs are bottom,
and at a boolean they select the same related branch. Application uses
the arrow clause; abstraction uses the body hypothesis for every related
argument pair. In recursion the bottom approximants are related, the
body hypothesis preserves related pairs, and directed closure relates
the least fixed points. Thus the binary lemma holds for all terms.
For any closing context of two closed terms, choose a new hole variable
different from every context binder. Its context term is typed with
that variable as its only free declaration. The binary lemma on the
two hole environments and semantic substitution give a ground pair in
\(B_g\), hence equality. The context's existing binders are retained.
No different domain interpretation's adequacy is used.
\end{proof}
\end{annexitem}

\begin{annexitem}{annex:p10.oa.not-fully-abstract}{thm:p10-oa-not-fully-abstract}{Sequential test impossibility and the semantic separator}
\begin{proof}
Let \(h\in A_{\mathsf{bool}\to\mathsf{bool}\to\mathsf{bool}}\).
Both \((\mathsf{true},\bot,\mathsf{false})\) and
\((\bot,\mathsf{true},\mathsf{false})\) belong to \(Q_{\mathsf{bool}}\).
Apply the lifted relation twice to \((h,h,h)\). It gives
\[
 (h(\mathsf{true})(\bot),\ h(\bot)(\mathsf{true}),\
                  h(\mathsf{false})(\mathsf{false}))\in Q_{\mathsf{bool}}.
\]
If the first two components are true, the third must be true. Each
of \(T,F\) therefore sends such an \(h\) to bottom: a false outer
guard selects \(\Omega\), a bottom outer guard produces bottom, and
when both outer guards are true the inner guard is true and again
selects \(\Omega\). This reasoning covers all possible guard values.
The terms' own denotations lie in the appropriate \(A\) by the closed
ternary lemma. Every binary-related input pair has sequential endpoints;
the two outputs are bottom, hence binary-related. Thus the denotations
of \(T,F\) are binary-related at their arrow type. The context lemma
gives equal ground denotations in every closing context, and the proved
ground adequacy yields the same value or divergence in each. Therefore
\(T\simeq_o F\).

The table for \(\operatorname{por}\) is monotone on the finite
flat-boolean product. A true coordinate keeps the output true under
refinement; false/false is maximal with output false; other nonmaximal
outputs are bottom, below each refined output. Every directed subset
of a finite poset has a greatest member, obtained by finitely many
directedness steps, so any monotone map from it is Scott-continuous.
Currying therefore places this table in the full continuous function
carrier. Its first two tested outputs are true and its third tested
output false. The innermost else-branch consequently returns true in
\(T\) and false in \(F\), proving unequal denotations.
The same triple test shows that por is not in \(A\), so its presence
in the full model does not contradict closed-term sequentiality.
This proves failure of full abstraction, not the stronger full-abstraction
claim for an extended language containing a parallel primitive.
\end{proof}
\end{annexitem}
